\chapter{Digital Signatures}

In Chapter~21, we saw how elliptic curves make Diffie--Hellman fast and compact, while
leaving open how either party verifies who is on the other end of the connection. In the
symmetric setting, we solved message authenticity with a MAC (Chapter~15), which proves
that a message came from \emph{somebody who holds the shared key}. That is enough when Alice and Bob already share a secret and trust each other. But it has two limits that turn out to matter enormously in the open world. First, the receiver can also \emph{produce} valid tags -- the same key both signs and verifies -- so Bob cannot prove to a third party that Alice, and not Bob himself, made the message. Second, before any of this works, Alice and Bob must already share a secret, which is exactly the bootstrapping problem public-key cryptography was invented to solve.

A \emph{digital signature} fixes both. Alice signs with a \emph{private} key that only she holds; anyone with her \emph{public} key can verify. The verifier cannot forge -- verification and signing use different keys -- so a valid signature is evidence, checkable by the whole world, that the holder of Alice's private key endorsed this exact message. That is \emph{public verifiability}, and it brings with it \emph{non-repudiation}: Alice cannot later credibly deny having signed, because nobody else could have produced the signature.

Signatures underpin most practical uses of public-key cryptography. Every time your browser trusts a certificate, every software update your phone installs, every block added to a cryptocurrency ledger, every SSH host key you accept -- a signature was verified. This chapter builds them from the ground up: what they promise and what they do \emph{not}; why ``textbook'' RSA signatures are dangerous and why we always hash-then-sign; the discrete-log signatures DSA, ECDSA, and EdDSA/Ed25519; and one of the most damaging implementation bugs in the field's history -- nonce reuse in ECDSA, the bug that broke code signing on the Sony PlayStation~3 and emptied Bitcoin wallets. Small numerical examples run throughout.

\chapterline

\section{Signatures Versus MACs}
\label{ch22:sig-vs-mac}

Both a MAC and a signature attach a short tag to a message so the receiver can detect tampering and confirm origin. The difference is \emph{who can make a valid tag} and \emph{who can check it}.

\begin{itemize}
  \item \textbf{MAC:} one secret key $k$ does everything. To make a tag you need $k$; to verify a tag you need $k$. Sender and verifier are symmetric -- either can forge against the other.
  \item \textbf{Signature:} a key \emph{pair} $(pk, sk)$. The private $sk$ makes signatures; the public $pk$ only verifies them. The signer and the verifier are asymmetric -- the verifier provably cannot forge.
\end{itemize}

This asymmetry between signer and verifier gives signatures three capabilities that symmetric MACs lack.

\begin{itemize}
  \item \textbf{Public verifiability.} Anyone holding $pk$ -- which is not secret and can be published on a billboard -- can check the signature. A MAC tag is checkable only by the holder of $k$.
  \item \textbf{No shared-secret setup.} The verifier needs only Alice's public key, which she can hand out freely. No secret channel is required before the first verified message.
  \item \textbf{Non-repudiation.} Because only Alice could have produced a signature that verifies under her $pk$, she cannot plausibly deny it later. A MAC has no such property: since Bob also holds $k$, a tag ``from Alice'' could just as easily have been forged by Bob, so it proves nothing to a judge.
\end{itemize}

\begin{intuitionbox}
A MAC is a wax seal pressed with a signet ring that \emph{both} correspondents own. If you find a sealed letter, you know one of the two ring-holders made it -- but you cannot tell which, and you cannot show anyone else, because you own the ring too and might have faked it.

A signature is more like a handwritten signature backed by a worldwide notary: Alice alone can produce her mark (the private key), but \emph{everyone} carries a reference copy to check it against (the public key). The asymmetry -- easy to verify, impossible to forge -- is what lets a stranger, or a court, rely on it.
\end{intuitionbox}

\begin{focusbox}
Non-repudiation is a feature \emph{and} a liability. Signatures are exactly the wrong tool for a casual chat where you might later want deniability. This is why Signal's messaging (Chapter~27) deliberately uses shared-key MACs for per-message authentication and reserves signatures for the long-term identity step -- so that a leaked transcript proves nothing about who said what. ``Deniable authentication'' is a real design goal, and signatures destroy it.
\end{focusbox}

\subsection{The syntax of a signature scheme}
\label{ch22:syntax}

\begin{definition}[Signature scheme]
\label{ch22:def-scheme}
A \emph{signature scheme} is a triple of algorithms $(\Gen, \Sign, \Vrfy)$:
\begin{itemize}
  \item $\Gen(1^n) \to (pk, sk)$: a randomized key generator producing a public verification key $pk$ and a private signing key $sk$.
  \item $\Sign_{sk}(m) \to \sigma$: takes the private key and a message, returns a signature $\sigma$. May be randomized.
  \item $\Vrfy_{pk}(m, \sigma) \to \{0,1\}$: takes the public key, a message, and a candidate signature, returns $1$ (accept) or $0$ (reject).
\end{itemize}
\emph{Correctness:} for every key pair from $\Gen$ and every message $m$, $\Vrfy_{pk}(m, \Sign_{sk}(m)) = 1$.
\end{definition}

\subsection{What ``secure'' means: existential unforgeability}
\label{ch22:euf-cma}

The security target mirrors the MAC definition, with one big change: the attacker is now \emph{given the public key}. She can verify on her own all day long; she does not even need an oracle to verify. What she must not be able to do is produce a \emph{new} valid message--signature pair.

\begin{formalbox}
\textbf{EUF-CMA (existential unforgeability under chosen-message attack).} Run $\Gen$ to get $(pk, sk)$ and hand $pk$ to the adversary $\Adv$. Give $\Adv$ a signing oracle: she submits messages $m_1, m_2, \dots$ of her choosing and receives $\sigma_i = \Sign_{sk}(m_i)$. Eventually $\Adv$ outputs a pair $(m^*, \sigma^*)$. She \emph{wins} if
\[
\Vrfy_{pk}(m^*, \sigma^*) = 1 \quad\text{and}\quad m^* \notin \{m_1, m_2, \dots\}.
\]
The scheme is EUF-CMA secure if every $\PPT$ adversary wins with only negligible probability:
\[
\Pr[\Adv \text{ wins}] \le \negl(n).
\]
\end{formalbox}

``Existential'' means the forgery can be on \emph{any} message $\Adv$ likes -- even meaningless gibberish counts as a break. ``Chosen-message'' means she may first collect signatures on messages she picks. This is a deliberately paranoid bar: we demand that even after seeing signatures on a million chosen documents, she cannot forge on a single fresh one. A scheme that clears EUF-CMA is safe against far weaker real-world attackers.

\begin{trapbox}
A subtle but important strengthening is \emph{strong} unforgeability (sEUF-CMA): $\Adv$ wins even if $m^*$ \emph{was} queried, as long as $\sigma^*$ is a \emph{new} signature on it. Plain EUF-CMA permits ``signature malleability'' -- producing a different valid signature on an already-signed message. That sounds harmless, but it has caused real damage: Bitcoin's pre-2017 ``transaction malleability'' bug let attackers alter a transaction's signature (and thus its ID) without invalidating it, which was used to confuse exchanges. We will see ECDSA's malleability concretely in Section~\ref{ch22:ecdsa}.
\end{trapbox}

\section{RSA Signatures}
\label{ch22:rsa-sig}

RSA (Chapter~19) gives us a trapdoor permutation on $\Zns{N}$: raising to the public exponent $e$ is easy for everyone, and the inverse -- raising to the private exponent $d$ -- is easy only for the key-holder. Encryption used $e$ to scramble and $d$ to unscramble. Signing flips the roles: the signer applies the \emph{private} operation, and verification applies the \emph{public} one. Anyone can run the public map forward; only the key-holder can run it backward; so only the key-holder can produce a value whose public image is what it should be.

\begin{historybox}
Diffie and Hellman named the goal --- a public-key signature whose forgery is as hard as
inverting a one-way function --- in their 1976 paper, and Rivest, Shamir, and Adleman's
1978 paper gave the first construction alongside RSA encryption: ``a message can be
\emph{signed} by deciphering it with the secret deciphering key.'' Discrete-log
signatures followed with Taher ElGamal (1985) and Claus Schnorr (1989). Because Schnorr's
signature was patent-encumbered until 2008, NIST's 1991 federal standard, the
\emph{Digital Signature Algorithm} (DSA, in FIPS~186), was deliberately designed to
sidestep that patent; it was also mildly controversial because some suspected NIST and
the NSA of pushing a signature-only scheme they understood especially well. Once
Schnorr's patent expired, a 2011 paper by Bernstein, Duif, Lange, Schwabe, and Yang
(``High-speed high-security signatures'') paired Schnorr's clean equation with
edwards25519 to create EdDSA and Ed25519 (RFC~8032), which won rapid adoption across SSH,
TLS~1.3, and Signal.
\end{historybox}

\subsection{Textbook RSA, and why it is broken}
\label{ch22:textbook-rsa}

The naive (``textbook'') scheme is the direct dual of textbook RSA encryption. Keys: $N = pq$, public exponent $e$, private exponent $d$ with $ed \equiv 1 \pmod{\varphi(N)}$.
\[
\Sign_{sk}(m) = m^{d} \bmod N, \qquad
\Vrfy_{pk}(m, \sigma): \text{ accept iff } \sigma^{e} \equiv m \pmod N.
\]
Correctness is the same Euler-theorem calculation that makes RSA decryption work: $\sigma^e = (m^d)^e = m^{ed} = m^{1 + k\varphi(N)} \equiv m \pmod N$.

Let us trace it with the textbook tiny primes.

\begin{example}[Textbook RSA signing by hand]
\label{ch22:ex-textbook}
Take $p = 61$, $q = 53$, so $N = 3233$ and $\varphi(N) = 60 \cdot 52 = 3120$. Pick $e = 17$; then $d = 2753$ since $17 \cdot 2753 = 46801 = 15\cdot 3120 + 1$.

Sign the ``message'' $m = 100$. The signature is
\[
\sigma = 100^{2753} \bmod 3233 = 1391.
\]
(Compute by square-and-multiply; the algorithms companion covers the cost.) To verify, the public side computes
\[
\sigma^{e} = 1391^{17} \bmod 3233 = 100 = m. \;\checkmark
\]
The verifier used only the public $(N, e) = (3233, 17)$ and recovered $m$, confirming the signer knew $d$.
\end{example}

The scheme is correct but badly insecure. Three attacks break it.

\begin{attackbox}
\textbf{1. Forge on a random ``message''.} The forger does not need to choose the message first. Pick any $\sigma$ you like, compute $m = \sigma^{e} \bmod N$, and announce ``$\sigma$ is my signature on $m$.'' It verifies by construction. The message $m$ is garbage you did not control -- but EUF-CMA says forging on \emph{any} message is a win, and in a system where messages carry meaning (a bank balance, a serial number) a garbage message can still be damaging or can be massaged into something useful.

\textbf{2. Multiplicative forgery (the deadly one).} Textbook RSA is \emph{homomorphic}: $(m_1 m_2)^d \equiv m_1^d \cdot m_2^d \pmod N$. So if Eve gets Alice to sign $m_1$ and $m_2$, she gets a free signature on $m_1 m_2 \bmod N$ that Alice never produced:
\[
\sigma_1 = m_1^d,\quad \sigma_2 = m_2^d \;\Longrightarrow\; \sigma_1\sigma_2 \bmod N = (m_1 m_2)^d = \Sign(m_1 m_2 \bmod N).
\]
Worse, to forge on a \emph{target} $m^*$, Eve writes $m^* = m_1 \cdot m_2 \bmod N$ for messages $m_1, m_2$ she chooses, asks Alice to sign those two (innocuous-looking) values, and multiplies the results. This is a textbook chosen-message break.

\textbf{3. No binding to one message.} Nothing ties $\sigma$ to a single $m$; the structure is pure algebra, so algebraic relations among messages carry straight over to signatures.
\end{attackbox}

\begin{example}[Multiplicative forgery with tiny numbers]
\label{ch22:ex-mult-forge}
Continue with $N = 3233$, $e = 17$, $d = 2753$. Eve wants a signature on $m^* = 4$ but Alice refuses to sign $4$ directly. Eve factors the target: $4 = 2 \cdot 2 \bmod 3233$. She asks Alice to sign $m_1 = 2$:
\[
\sigma_1 = 2^{2753} \bmod 3233 = 1027.
\]
Then $\sigma^* = \sigma_1^2 \bmod 3233 = 1027^2 \bmod 3233 = 1054729 \bmod 3233 = 771$. Check: $771^{17} \bmod 3233 = 4 = m^*$. Eve forged a signature on $4$ from a signature on $2$. No knowledge of $d$ required.
\end{example}

\subsection{Hash-then-sign: the fix}
\label{ch22:hash-then-sign}

Every attack above exploited \emph{algebraic structure} in the message -- the forger picked $m$'s that multiply nicely or read $m$ off backward from $\sigma$. The fix is to destroy that structure before signing: run the message through a cryptographic hash $\Hash$ first, then sign the digest (Figure~\ref{ch22:fig-rsa-hash-vs-textbook}).
\[
\Sign_{sk}(m) = \Hash(m)^{d} \bmod N, \qquad
\Vrfy_{pk}(m, \sigma): \text{ accept iff } \sigma^{e} \equiv \Hash(m) \pmod N.
\]

\begin{figure}[H]
\centering
\begin{tikzpicture}[font=\small, >=stealth]
  % Top panel: Textbook RSA multiplicative forgery
  \node[anchor=south west, font=\footnotesize\scshape, text=trapred!90!black] at (0,3.9) {Textbook RSA: Homomorphism Enables Forgery};
  \draw[draw=trapred!40, fill=trapred!5, rounded corners=2pt] (0,2.2) rectangle (14.5,3.8);

  \node[draw=aliceblue!60, fill=aliceblue!12, rounded corners=2pt, font=\scriptsize, inner sep=4pt] (t1) at (2.1,3.0) {Queries $m_1, m_2$};
  \node[draw=aliceblue!60, fill=aliceblue!18, rounded corners=2pt, font=\scriptsize, inner sep=4pt] (t2) at (6.0,3.0) {$\sigma_1 = m_1^d,\;\sigma_2 = m_2^d$};
  \node[draw=trapred!70, fill=trapred!15, rounded corners=2pt, font=\scriptsize\bfseries, inner sep=4pt] (t3) at (11.2,3.0) {Forge $\sigma^* = \sigma_1\sigma_2 \equiv (m_1 m_2)^d \pmod N$};

  \draw[->, thick, aliceblue!80!black] (t1) -- node[above, font=\tiny]{Oracle} (t2);
  \draw[->, thick, trapred!80!black] (t2) -- node[above, font=\tiny]{Multiply} (t3);

  % Bottom panel: Hash-then-sign breaks homomorphism
  \node[anchor=south west, font=\footnotesize\scshape, text=targetgreen!85!black] at (0,1.6) {Hash-Then-Sign RSA: One-Way Hash Breaks Algebraic Structure};
  \draw[draw=targetgreen!40, fill=targetgreen!5, rounded corners=2pt] (0,-0.1) rectangle (14.5,1.5);

  \node[draw=sourceblue!60, fill=sourceblue!12, rounded corners=2pt, font=\scriptsize, inner sep=4pt] (h1) at (1.5,0.7) {Message $m$};
  \node[draw=structurepurple!60, fill=structurepurple!12, rounded corners=2pt, font=\scriptsize, inner sep=4pt] (h2) at (4.5,0.7) {Digest $h = \Hash(m)$};
  \node[draw=targetgreen!70, fill=targetgreen!15, rounded corners=2pt, font=\scriptsize, inner sep=4pt] (h3) at (8.1,0.7) {$\sigma = \Hash(m)^d \bmod N$};
  \node[draw=gray!60, fill=white, rounded corners=2pt, font=\scriptsize, align=center, inner sep=4pt] (h4) at (12.1,0.7) {$\Hash(m_1 m_2) \neq \Hash(m_1)\Hash(m_2)$\\[-1pt](no multiplicative shortcut)};

  \draw[->, thick, gray!70] (h1) -- (h2);
  \draw[->, thick, gray!70] (h2) -- (h3);
  \draw[->, thick, targetgreen!80!black, dashed] (h3) -- (h4);
\end{tikzpicture}
\caption{Why we always hash before signing. Top: In textbook RSA, the signing map $m \mapsto m^d \bmod N$ is multiplicatively homomorphic, allowing an attacker to multiply two valid signatures into a forgery on $m_1 m_2$. Bottom: Applying a cryptographic hash $\Hash$ before raising to $d$ destroys the multiplicative relation between messages and digests.}
\label{ch22:fig-rsa-hash-vs-textbook}
\end{figure}

\begin{whybox}
Hashing breaks all three attacks at once.
\begin{itemize}
  \item \textbf{Against the random-$\sigma$ forgery:} given a random $\sigma$, the forger gets a random target $t = \sigma^e \bmod N$. To pass verification she now needs a message $m$ with $\Hash(m) = t$. That is a \emph{preimage} on the hash -- infeasible for SHA-256. She can pick $\sigma$ or pick $m$, but not both consistently.
  \item \textbf{Against multiplicative forgery:} the homomorphism still holds at the level of integers, but now the forger needs three messages $m_1, m_2, m^*$ with $\Hash(m^*) \equiv \Hash(m_1)\cdot\Hash(m_2) \pmod N$. Hash outputs are effectively random in $\Zns{N}$; finding such a multiplicative relation among hash digests is as hard as a fresh search and provides no leverage.
  \item \textbf{Bonus -- arbitrary length:} RSA can only sign numbers below $N$. Hashing first compresses any message, however long, to a fixed-size digest.
\end{itemize}
\end{whybox}

\begin{trapbox}
Hash-then-sign with a \emph{bare} hash padded with zeros is still subtle. The deployed standard, \textbf{RSASSA-PKCS1-v1\_5}, sandwiches the digest in a rigid byte structure \texttt{00 01 FF FF ... FF 00 || ASN.1(hashID) || digest}; \textbf{RSA-PSS} adds randomized salt and a mask, with a tight security proof in the random-oracle model. The infamous \emph{Bleichenbacher signature-forgery} attack (2006) hit careless implementations that parsed the PKCS\#1 padding loosely: if a verifier did not check that the \texttt{FF} padding ran all the way to the right edge, an attacker could, for small $e=3$, construct a perfect cube whose low bits encoded a valid-\emph{looking} structure, forging without the private key. The lesson: verification must check the padding \emph{exactly}, byte for byte. Prefer PSS for new designs.
\end{trapbox}

\begin{applicationbox}
RSA signatures with $e = 65537$ remain ubiquitous. The majority of TLS certificates in the field (Chapter~24) are RSA-signed; code-signing for Windows drivers and many firmware images uses RSA; DNSSEC, JWT (\texttt{RS256}), and PDF signing all rely on it. New systems increasingly prefer the smaller, faster Ed25519 (Section~\ref{ch22:eddsa}), but RSA's two-decade install base means it will be verified for a long time yet. When you see \texttt{RSASSA-PSS} or \texttt{sha256WithRSAEncryption} in a certificate, this is the machinery behind it.
\end{applicationbox}

\begin{implementbox}
A toy RSA signature, hash-then-sign, in Python. \textbf{For learning only -- never use this for anything real.} It uses tiny keys, a truncated hash reduced mod a small $N$, no PKCS\#1/PSS padding, and Python's non-constant-time big-int arithmetic. Every one of those is a fatal flaw in production.
\begin{lstlisting}[language=Python]
import hashlib

# Toy keypair (textbook tiny primes -- NOT secure)
p, q = 61, 53
N = p * q                      # 3233
phi = (p - 1) * (q - 1)        # 3120
e = 17
d = pow(e, -1, phi)            # 2753

def H(msg: bytes) -> int:
    # Hash then reduce into Z_N. Real RSA uses full-width PSS padding.
    digest = hashlib.sha256(msg).digest()
    return int.from_bytes(digest, "big") % N

def sign(msg: bytes) -> int:
    return pow(H(msg), d, N)              # private op: ^d

def verify(msg: bytes, sigma: int) -> bool:
    return pow(sigma, e, N) == H(msg)     # public op: ^e

s = sign(b"Pay Bob 100")
print(verify(b"Pay Bob 100", s))          # True
print(verify(b"Pay Eve 9999", s))         # False -- different hash
\end{lstlisting}
Tampering with the message changes the digest, so the recomputed $\Hash(m)$ no longer matches $\sigma^e$, and verification fails. These few lines show integrity and authenticity at work, but the warnings above still apply.
\end{implementbox}

\section{Discrete-Log Signatures: DSA}
\label{ch22:dsa}

RSA signatures rest on the hardness of factoring. A second family rests on
\emph{discrete logarithms} (Chapter~20): given $g$ and $g^x$ in a cyclic group, recovering
$x$ is hard. The Digital Signature Algorithm (DSA), standardized by NIST in 1991, brought
discrete-log signatures into federal standards and established the foundation for modern
elliptic-curve variants. Because its signatures are substantially shorter than RSA's,
the structure is worth understanding in detail.

\subsection{Setup and the algorithm}
\label{ch22:dsa-algo}

DSA works in a subgroup of $\Zns{p}$. Public parameters: a large prime $p$, a prime $q$ dividing $p-1$, and a generator $g$ of the order-$q$ subgroup (so $g^q \equiv 1 \pmod p$). The private key is $x \getsr \{1,\dots,q-1\}$; the public key is $y = g^{x} \bmod p$.

To sign a message $m$ with digest $h = \Hash(m) \bmod q$:
\begin{enumerate}
  \item Choose a fresh random \emph{nonce} $k \getsr \{1, \dots, q-1\}$.
  \item Compute $r = (g^{k} \bmod p) \bmod q$. If $r = 0$, restart with a new $k$.
  \item Compute $s = k^{-1}(h + x r) \bmod q$. If $s = 0$, restart.
  \item The signature is $(r, s)$.
\end{enumerate}

To verify $(r,s)$ on $m$ (with $h = \Hash(m) \bmod q$), reject unless $0 < r < q$ and $0 < s < q$, then:
\[
w = s^{-1} \bmod q,\quad u_1 = h\,w \bmod q,\quad u_2 = r\,w \bmod q,\quad v = \big(g^{u_1} y^{u_2} \bmod p\big)\bmod q.
\]
Accept iff $v = r$.

\begin{whybox}
\textbf{Why verification works.} From step~3, $s = k^{-1}(h + xr) \bmod q$, so $k \equiv s^{-1}(h + xr) \equiv w(h + xr) \equiv wh + wxr \equiv u_1 + x u_2 \pmod q$. Now raise $g$ to both sides, working in the order-$q$ subgroup where exponents only matter mod $q$:
\[
g^{k} \equiv g^{u_1 + x u_2} = g^{u_1}(g^{x})^{u_2} = g^{u_1} y^{u_2} \pmod p.
\]
Reducing both sides mod $q$ gives $r = v$. The signer's secret $x$ appears on the right only through the public $y = g^x$, so the verifier never needs $x$ -- exactly the asymmetry we want.
\end{whybox}

\begin{example}[DSA by hand with tiny parameters]
\label{ch22:ex-dsa}
Let $p = 23$, $q = 11$ (note $11 \mid 22$). We need $g$ of order $11$. Take $h_0 = 2$ and set $g = h_0^{(p-1)/q} = 2^{2} = 4 \bmod 23$. Check: $4^{11} \bmod 23 = 1$, and $4 \neq 1$, so $\ord(g) = 11$. Good.

Private key $x = 3$; public key $y = g^x = 4^3 = 64 \bmod 23 = 18$.

Sign a message whose digest is $h = 9$. Pick nonce $k = 7$ (with $k^{-1} \bmod 11 = 8$, since $7\cdot 8 = 56 = 5\cdot 11 + 1$).
\begin{align*}
r &= (g^{k} \bmod p)\bmod q = (4^{7} \bmod 23)\bmod 11.\\
4^7 &= 16384,\quad 16384 \bmod 23 = 8,\quad r = 8 \bmod 11 = 8.\\
s &= k^{-1}(h + x r)\bmod q = 8\,(9 + 3\cdot 8)\bmod 11 = 8\cdot 33 \bmod 11 = 8\cdot 0 = 0.
\end{align*}
Here $33 \equiv 0 \pmod{11}$ gives $s = 0$, which DSA forbids -- so we restart. Take $k = 4$ instead (with $k^{-1} \bmod 11 = 3$, since $4\cdot 3 = 12 \equiv 1$).
\begin{align*}
r &= (4^{4}\bmod 23)\bmod 11 = (256 \bmod 23)\bmod 11 = 3 \bmod 11 = 3.\\
s &= 3\,(9 + 3\cdot 3)\bmod 11 = 3\cdot 18 \bmod 11 = 54 \bmod 11 = 10.
\end{align*}
Signature: $(r, s) = (3, 10)$.

\textbf{Verify.} $w = s^{-1} = 10^{-1}\bmod 11 = 10$ (since $10\cdot 10 = 100 = 9\cdot 11 + 1$).
\begin{align*}
u_1 &= h w = 9\cdot 10 = 90 \bmod 11 = 2,\\
u_2 &= r w = 3\cdot 10 = 30 \bmod 11 = 8,\\
v &= (g^{u_1} y^{u_2}\bmod 23)\bmod 11 = (4^{2}\cdot 18^{8} \bmod 23)\bmod 11.
\end{align*}
Now $4^2 = 16$. And $18^2 = 324 \bmod 23 = 2$, so $18^4 = 2^2 = 4$, $18^8 = 4^2 = 16 \bmod 23$. Thus $g^{u_1}y^{u_2} = 16\cdot 16 = 256 \bmod 23 = 3$, and $v = 3 \bmod 11 = 3 = r$. Accept. \checkmark
\end{example}

\begin{trapbox}
DSA's security depends directly on how the nonce $k$ is generated, and flawed nonce generation is the most common way real implementations fail. It must be \emph{uniformly random}, \emph{secret}, and \emph{unique per signature}. Reusing it, or letting it be predictable or biased, leaks the private key outright. Section~\ref{ch22:nonce-attack} is devoted to exactly this, because nonce failures have broken real, shipped systems.
\end{trapbox}

\section{ECDSA: DSA on an Elliptic Curve}
\label{ch22:ecdsa}

ECDSA is DSA with the multiplicative group $\Zns{p}$ replaced by the additive group of points on an elliptic curve (Chapter~21). The benefit is the usual one for elliptic curves: equivalent security at far smaller key and signature sizes, because the best discrete-log attack on a well-chosen curve is the generic $\sqrt{|\text{group}|}$ square-root attack. A 256-bit curve gives roughly 128-bit security; matching that with DSA over $\Zns{p}$ needs a 3072-bit $p$.

Public parameters: a curve $E$ over $\F_p$, a base point $G$ of large prime order $n$. Private key $d \getsr \{1,\dots,n-1\}$; public key $Q = dG$ (scalar multiplication, the elliptic-curve analogue of $g^x$).

\textbf{Sign} $m$ with $h = \Hash(m) \bmod n$:
\begin{enumerate}
  \item Pick a fresh random nonce $k \getsr \{1,\dots,n-1\}$.
  \item Compute the point $(x_1, y_1) = kG$; set $r = x_1 \bmod n$. If $r = 0$, restart.
  \item Compute $s = k^{-1}(h + d\,r) \bmod n$. If $s = 0$, restart.
  \item Signature is $(r, s)$.
\end{enumerate}

\textbf{Verify} $(r,s)$ on $m$: reject unless $0 < r,s < n$; then
\[
w = s^{-1}\bmod n,\quad u_1 = hw \bmod n,\quad u_2 = rw \bmod n,\quad (x_1,y_1) = u_1 G + u_2 Q.
\]
Accept iff $x_1 \bmod n = r$.

The correctness argument is identical in shape to DSA's: $k \equiv w(h + dr) \equiv u_1 + d u_2 \pmod n$, so $kG = u_1 G + d u_2 G = u_1 G + u_2 Q$, and the $x$-coordinates match.

\begin{figure}[H]
\centering
\begin{tikzpicture}[font=\small, >=stealth]
  % Lane 1: Signing (Left)
  \begin{scope}[shift={(0,0)}]
    \node[anchor=south west, font=\footnotesize\scshape, text=sourceblue!80!black] at (0,7.2) {Signing Pipeline (Signer holding private key $d$)};
    \draw[draw=sourceblue!40, fill=sourceblue!5, rounded corners=2pt] (0,0) rectangle (6.5,7.0);

    \node[draw=sourceblue!60, fill=sourceblue!15, rounded corners=2pt, font=\scriptsize, align=left, text width=5.8cm] (s_in) at (3.25,6.3) {
      \textbf{Inputs}: Message $m$, private key $d$, generator $G$, curve order $n$
    };

    \node[draw=gray!50, fill=white, rounded corners=2pt, font=\scriptsize, align=left, text width=5.8cm] (s_hash) at (3.25,5.1) {
      1. Hash message: $h = \Hash(m) \bmod n$
    };

    \node[draw=quotientorange!70, fill=quotientorange!10, rounded corners=2pt, font=\scriptsize, align=left, text width=5.8cm] (s_nonce) at (3.25,3.9) {
      2. Choose fresh ephemeral secret nonce: $k \stackrel{\$}{\gets} \{1, \dots, n-1\}$
    };

    \node[draw=gray!50, fill=white, rounded corners=2pt, font=\scriptsize, align=left, text width=5.8cm] (s_point) at (3.25,2.7) {
      3. Compute point: $(x_1, y_1) = k G$\\
      \phantom{3. }Set $r = x_1 \bmod n$ (restart if $r=0$)
    };

    \node[draw=gray!50, fill=white, rounded corners=2pt, font=\scriptsize, align=left, text width=5.8cm] (s_eq) at (3.25,1.5) {
      4. Bind scalar: $s = k^{-1}(h + d\,r) \bmod n$ (restart if $s=0$)
    };

    \node[draw=sourceblue!80, fill=sourceblue!20, rounded corners=2pt, font=\scriptsize\bfseries, align=center, text width=5.8cm] (s_out) at (3.25,0.5) {
      Output Signature: $\sigma = (r, s)$
    };

    \draw[->, thick, gray!70] (s_in) -- (s_hash);
    \draw[->, thick, gray!70] (s_hash) -- (s_nonce);
    \draw[->, thick, gray!70] (s_nonce) -- (s_point);
    \draw[->, thick, gray!70] (s_point) -- (s_eq);
    \draw[->, thick, gray!70] (s_eq) -- (s_out);
  \end{scope}

  % Public Transmission Corridor
  \draw[->, very thick, quotientorange!80!black] (6.6,0.5) -- node[above, font=\scriptsize, text=black]{$\sigma = (r,s), \; m$} (7.9,0.5);

  % Lane 2: Verification (Right)
  \begin{scope}[shift={(8.0,0)}]
    \node[anchor=south west, font=\footnotesize\scshape, text=targetgreen!80!black] at (0,7.2) {Verification Pipeline (Public using $Q = dG$)};
    \draw[draw=targetgreen!40, fill=targetgreen!5, rounded corners=2pt] (0,0) rectangle (6.5,7.0);

    \node[draw=targetgreen!60, fill=targetgreen!15, rounded corners=2pt, font=\scriptsize, align=left, text width=5.8cm] (v_in) at (3.25,6.3) {
      \textbf{Inputs}: Message $m$, signature $(r,s)$, public key $Q = dG$
    };

    \node[draw=gray!50, fill=white, rounded corners=2pt, font=\scriptsize, align=left, text width=5.8cm] (v_range) at (3.25,5.1) {
      1. Check ranges: $1 \le r, s < n$, and hash $h = \Hash(m) \bmod n$
    };

    \node[draw=gray!50, fill=white, rounded corners=2pt, font=\scriptsize, align=left, text width=5.8cm] (v_inv) at (3.25,3.9) {
      2. Invert signature scalar: $w = s^{-1} \bmod n$
    };

    \node[draw=gray!50, fill=white, rounded corners=2pt, font=\scriptsize, align=left, text width=5.8cm] (v_weights) at (3.25,2.7) {
      3. Compute linear weights:\\
      \phantom{3. }$u_1 = h \cdot w \bmod n, \quad u_2 = r \cdot w \bmod n$
    };

    \node[draw=gray!50, fill=white, rounded corners=2pt, font=\scriptsize, align=left, text width=5.8cm] (v_recon) at (3.25,1.5) {
      4. Reconstruct curve point:\\
      \phantom{4. }$(x_1', y_1') = u_1 G + u_2 Q$
    };

    \node[draw=targetgreen!80, fill=targetgreen!20, rounded corners=2pt, font=\scriptsize\bfseries, align=center, text width=5.8cm] (v_out) at (3.25,0.5) {
      Accept if $x_1' \bmod n = r$; otherwise Reject
    };

    \draw[->, thick, gray!70] (v_in) -- (v_range);
    \draw[->, thick, gray!70] (v_range) -- (v_inv);
    \draw[->, thick, gray!70] (v_inv) -- (v_weights);
    \draw[->, thick, gray!70] (v_weights) -- (v_recon);
    \draw[->, thick, gray!70] (v_recon) -- (v_out);
  \end{scope}
\end{tikzpicture}
\caption{The ECDSA signature generation and verification pipeline. The signer binds the
message hash $h$ and private key $d$ to an ephemeral point $kG$. The verifier uses
the public key $Q = dG$ to reconstruct the ephemeral point and confirms its $x$-coordinate matches $r$.}
\label{ch22:fig-ecdsa-pipeline}
\end{figure}

\begin{example}[ECDSA on a tiny curve]
\label{ch22:ex-ecdsa}
Use the curve $E: y^2 = x^3 + 2x + 2$ over $\F_{17}$, a standard small teaching curve. The point $G = (5, 1)$ has order $n = 19$. (You can verify $19 G = \mathcal{O}$; the group has $19$ points.)

Take private key $d = 7$. Then $Q = 7G$. Doubling and adding on this curve (Chapter~21's point arithmetic) gives the multiples; in particular $7G = (0, 6)$, so $Q = (0,6)$.

Sign a message with digest $h = 11$. Pick nonce $k = 3$. Compute $kG = 3G$. On this curve $3G = (10, 6)$, so $r = x_1 \bmod n = 10 \bmod 19 = 10$.

Now $k^{-1} \bmod 19$: $3\cdot 13 = 39 = 2\cdot 19 + 1$, so $k^{-1} = 13$.
\[
s = k^{-1}(h + d r)\bmod n = 13\,(11 + 7\cdot 10)\bmod 19 = 13\cdot 81 \bmod 19.
\]
$81 \bmod 19 = 81 - 76 = 5$, so $s = 13\cdot 5 = 65 \bmod 19 = 65 - 57 = 8$. Signature $(r,s) = (10, 8)$.

\textbf{Verify.} $w = s^{-1} = 8^{-1}\bmod 19 = 12$ (since $8\cdot 12 = 96 = 5\cdot 19 + 1$).
\[
u_1 = hw = 11\cdot 12 = 132 \bmod 19 = 132 - 114 = 18,\quad u_2 = rw = 10\cdot 12 = 120 \bmod 19 = 6.
\]
Compute $u_1 G + u_2 Q = 18G + 6Q$. Since $Q = 7G$, this is $18G + 42G = 60G$. And $60 \bmod 19 = 60 - 57 = 3$, so the result is $3G = (10,6)$, with $x_1 = 10$. Then $x_1 \bmod n = 10 = r$. Accept. \checkmark

The shortcut $Q = 7G \Rightarrow u_2 Q = 42G$ is just for hand-checking; a real verifier does the two scalar multiplications directly without knowing $d$.
\end{example}

\begin{trapbox}
\textbf{ECDSA is malleable.} If $(r, s)$ is a valid signature, so is $(r,\, n - s)$: negating $s$ corresponds to negating the nonce point's $y$-coordinate, which leaves $r = x_1 \bmod n$ unchanged. Both verify. This is the source of Bitcoin's transaction-malleability headaches; the fix (BIP~62 / ``low-$s$'' rule) is to canonicalize by requiring $s \le n/2$, rejecting the high-$s$ twin. EdDSA avoids the whole issue by being deterministic and strongly unforgeable.
\end{trapbox}

\begin{applicationbox}
ECDSA over the NIST curve P-256 (and increasingly P-384) is the default signature in modern TLS certificates and in the TLS handshake's \texttt{CertificateVerify}. It signs Apple and Android app packages, secures Bluetooth pairing, and -- over the Koblitz curve secp256k1 -- authorizes every Bitcoin and Ethereum transaction. Its compactness (a 64-byte signature for 128-bit security) is why it displaced RSA in size-constrained settings. Its sharp edge is the nonce, which the next section demonstrates concretely.
\end{applicationbox}

\section{The Catastrophic Nonce Attack}
\label{ch22:nonce-attack}

DSA and ECDSA share a fragility so severe it deserves its own section: the security of the \emph{long-term private key} rests entirely on the per-signature nonce $k$ being random, secret, and never repeated. Violate any of those and the key falls out with grade-school algebra. We show it for ECDSA; DSA is identical with $G$ replaced by $g$ and $n$ by $q$.

\subsection{Reusing one nonce reveals the key}
\label{ch22:nonce-reuse}

Suppose Alice signs two \emph{different} messages with digests $h_1 \neq h_2$ using the \emph{same} nonce $k$. Then both signatures share the same $r$, because $r$ depends only on $k$: $r = x\text{-coord}(kG) \bmod n$. The two $s$ values are
\[
s_1 = k^{-1}(h_1 + d r)\bmod n,\qquad s_2 = k^{-1}(h_2 + d r)\bmod n.
\]
Subtract. The $dr$ term cancels:
\[
s_1 - s_2 \equiv k^{-1}(h_1 - h_2) \pmod n \;\Longrightarrow\; k \equiv \frac{h_1 - h_2}{s_1 - s_2} \pmod n.
\]
That recovers the nonce. With $k$ in hand, rearrange either signature equation for the private key:
\[
s_1 k \equiv h_1 + d r \pmod n \;\Longrightarrow\; d \equiv \frac{s_1 k - h_1}{r} \pmod n.
\]
Two signatures, a subtraction, and two modular inverses recover the entire private key from public data. An attacker who merely \emph{observes} two same-nonce signatures (they are public!) computes $d$ and can forge anything forever.

\begin{figure}[H]
\centering
\begin{tikzpicture}[font=\small, >=stealth]
  % Lane 1: Two signatures with identical nonce
  \node[anchor=south west, font=\footnotesize\scshape, text=trapred!90!black] at (0,5.6) {Step 1: Public Observation of Two Signatures Sharing $r$};
  \draw[draw=trapred!40, fill=trapred!5, rounded corners=2pt] (0,3.6) rectangle (14.5,5.5);

  \node[draw=aliceblue!60, fill=aliceblue!15, rounded corners=2pt, font=\scriptsize, align=left, text width=6.6cm] at (3.6,4.55) {
    \textbf{Signature 1 on message $m_1$}: $({\color{trapred}\mathbf{r}}, s_1)$\\
    $s_1 \equiv k^{-1}(h_1 + d\,{\color{trapred}\mathbf{r}}) \pmod n \quad (h_1 = \Hash(m_1))$
  };

  \node[draw=bobgreen!60, fill=bobgreen!15, rounded corners=2pt, font=\scriptsize, align=left, text width=6.6cm] at (10.9,4.55) {
    \textbf{Signature 2 on message $m_2$}: $({\color{trapred}\mathbf{r}}, s_2)$\\
    $s_2 \equiv k^{-1}(h_2 + d\,{\color{trapred}\mathbf{r}}) \pmod n \quad (h_2 = \Hash(m_2))$
  };

  % Down arrow corridor
  \draw[->, very thick, trapred] (7.25,3.6) -- node[right, font=\scriptsize, text=trapred!90!black]{Identical $r \implies$ identical nonce $k = k_1 = k_2$} (7.25,2.7);

  % Lane 2: Canceling out the private key d
  \node[anchor=south west, font=\footnotesize\scshape, text=quotientorange!90!black] at (0,2.7) {Step 2: Subtract Signatures to Cancel Private Key $d$};
  \draw[draw=quotientorange!40, fill=quotientorange!5, rounded corners=2pt] (0,1.2) rectangle (14.5,2.6);

  \node[font=\scriptsize, align=center] at (7.25,1.9) {
    $s_1 - s_2 \equiv k^{-1}\big((h_1 + d\,r) - (h_2 + d\,r)\big) \equiv k^{-1}(h_1 - h_2) \pmod n$\\[0.3em]
    \textbf{Algebraic Nonce Extraction}: \quad
    ${\color{quotientorange!90!black}\mathbf{k \equiv (h_1 - h_2)(s_1 - s_2)^{-1} \pmod n}}$
  };

  % Down arrow corridor
  \draw[->, very thick, quotientorange!80!black] (7.25,1.2) -- node[right, font=\scriptsize, text=black]{Nonce $k$ is now completely known!} (7.25,0.4);

  % Lane 3: Private key extraction
  \node[anchor=south west, font=\footnotesize\scshape, text=evered!90!black] at (0,0.4) {Step 3: Solve for Long-Term Master Private Key $d$};
  \draw[draw=evered!50, fill=evered!10, rounded corners=2pt] (0,-1.0) rectangle (14.5,0.3);

  \node[font=\scriptsize, align=center] at (7.25,-0.35) {
    Rearrange signature equation: $s_1 k \equiv h_1 + d\,r \pmod n$\\[0.3em]
    \textbf{Complete Key Compromise}: \quad
    ${\color{evered}\mathbf{d \equiv (s_1 k - h_1)\,r^{-1} \pmod n}}$
  };
\end{tikzpicture}
\caption{The ECDSA repeated-nonce key extraction attack (the Sony PlayStation~3 disaster).
Because $r = (kG)_x$ depends solely on the nonce $k$, reusing $k$ across two distinct
messages produces the same $r$. Subtracting the signature equations eliminates the secret key $d$,
allowing the attacker to solve directly for $k$, and then immediately extract $d$.}
\label{ch22:fig-nonce-reuse}
\end{figure}

\begin{example}[Extracting the key from a reused nonce]
\label{ch22:ex-reuse}
Use the tiny ECDSA setup of Example~\ref{ch22:ex-ecdsa}: curve over $\F_{17}$, $n = 19$, true secret $d = 7$ (which the attacker does not know). Alice carelessly reuses $k = 3$ (so $r = 10$ in both) on two messages:
\[
h_1 = 11:\; s_1 = 8 \ (\text{from Example~\ref{ch22:ex-ecdsa}}),\qquad
h_2 = 5:\; s_2 = k^{-1}(h_2 + dr) = 13(5 + 70)\bmod 19.
\]
$5 + 70 = 75 \equiv 75 - 57 = 18 \pmod{19}$, so $s_2 = 13\cdot 18 = 234 \bmod 19 = 234 - 228 = 6$.

Eve sees $(r, s_1) = (10, 8)$ and $(r, s_2) = (10, 6)$ with the same $r$ -- the red flag. She recovers:
\[
k \equiv \frac{h_1 - h_2}{s_1 - s_2} = \frac{11 - 5}{8 - 6} = \frac{6}{2} = 3 \pmod{19}.
\]
($2^{-1}\bmod 19 = 10$, so $6\cdot 10 = 60 \bmod 19 = 3$. Indeed $k = 3$.) Then
\[
d \equiv \frac{s_1 k - h_1}{r} = \frac{8\cdot 3 - 11}{10} = \frac{13}{10}\bmod 19.
\]
$10^{-1}\bmod 19 = 2$ (since $10\cdot 2 = 20 \equiv 1$), so $d = 13\cdot 2 = 26 \bmod 19 = 7$. Eve has recovered the secret key $d = 7$ exactly, from two public signatures.
\end{example}

\begin{attackbox}
\textbf{The PlayStation~3 break (2010).} Sony signed PS3 firmware and games with ECDSA. The whole chain of trust rested on those signatures: only Sony-signed code would run. At the Chaos Communication Congress (27C3), the group \texttt{fail0verflow} revealed that Sony's signing implementation used a \emph{constant} nonce $k$ -- the same value for every signature, rather than a fresh random one each time. Every pair of signed binaries therefore shared $r$, and the two-signature subtraction above handed over Sony's master private key.

With the private key, anyone could sign arbitrary code that the PS3 would accept as genuine Sony software -- a complete, unrecoverable break of the platform's root of trust. There was no patch: you cannot un-leak a private key, and the public key was burned into millions of consoles. It is the standard real-world example of how ECDSA's security depends entirely on generating the nonce correctly.
\end{attackbox}

\begin{attackbox}
\textbf{Android Bitcoin wallet theft (2013).} A flaw in the \texttt{SecureRandom} of certain Android versions caused some Bitcoin wallet apps to generate \emph{repeated} ECDSA nonces across transactions. Because every Bitcoin signature is published forever on the blockchain, anyone scanning the chain for two signatures sharing an $r$ value could apply Example~\ref{ch22:ex-reuse}'s arithmetic, extract the wallet's private key, and sweep its funds. Real coins were stolen. The combination of permanently public signatures and a faulty RNG is especially dangerous: anyone can search the entire history for repeated nonces, and stolen coins cannot be recovered.
\end{attackbox}

\begin{connectionbox}
You do not even need a full \emph{repeat} to lose the key -- a slight \emph{bias} is enough. If each nonce leaks just a few bits (say its top byte is always zero, or it comes from a too-short range), then each signature gives a noisy linear equation in $d$, and gathering a few hundred such equations turns key recovery into a \emph{lattice} problem (the hidden number problem) that LLL-style lattice reduction solves. The probability companion's treatment of accumulating a faint statistical signal across many samples applies directly: each biased nonce is a weak measurement, and enough weak measurements pin down the secret. ``Random enough'' is therefore a precise, quantitative requirement.
\end{connectionbox}

\subsection{The fix: deterministic nonces (RFC 6979)}
\label{ch22:rfc6979}

The cleanest defense removes randomness from the nonce entirely. Instead of sampling $k$ from an RNG you may not trust, \emph{derive} it deterministically from the private key and the message:
\[
k = \mathrm{HMAC\text{-}derived}(d, \Hash(m)).
\]
This is RFC~6979. Because $k$ depends on the message, two different messages get different (and unpredictable, since $d$ is secret) nonces, so reuse is structurally impossible. Because it depends on $d$, an outsider cannot predict it. And because it needs no RNG at signing time, a broken \texttt{SecureRandom} can no longer sink you. The same private key and message always yield the same signature -- which is fine, and even desirable for reproducibility and testing.

\begin{whybox}
Determinism sounds paradoxical -- did we not say the nonce must be ``random''? The real requirement is weaker and more precise: each nonce must be \emph{unpredictable to the attacker} and \emph{distinct across distinct messages}. A keyed PRF (HMAC) of the message satisfies both: distinct messages give distinct, pseudorandom outputs, and without the key the outputs look uniform. Determinism just means we obtain that unpredictability from a hash of secret data rather than from a hardware RNG -- trading a fragile dependency for a robust one. EdDSA, next, bakes this in by design.
\end{whybox}

\section{EdDSA and Ed25519}
\label{ch22:eddsa}

EdDSA (Edwards-curve Digital Signature Algorithm), and its instantiation \emph{Ed25519},
is the modern signature you should reach for first. It pairs Schnorr's clean signature
equation with Bernstein's edwards25519 curve, standardized in RFC~8032. The design
philosophy is strict misuse-resistance: no per-signature randomness to mishandle, no
point-validation footguns, and complete, constant-time curve arithmetic. Today, SSH,
TLS~1.3, Signal, and modern operating systems use it as their default.

\subsection{How Ed25519 signs}
\label{ch22:ed25519-algo}

Parameters: the twisted Edwards curve edwards25519 over $\F_{2^{255}-19}$, base point $B$ of prime order $\ell \approx 2^{252}$, and the hash $H = $ SHA-512.

\textbf{Watch the letters.} We follow RFC~8032's names, which swap roles with ECDSA: here $r$ is the \emph{secret} nonce scalar and $k$ is the \emph{public} hash challenge, whereas in ECDSA above $k$ was the secret nonce and $r$ a public coordinate.

\textbf{Keys.} The private key is a 32-byte seed. Hash it with SHA-512 to get 64 bytes; the first 32 (clamped) form the scalar $a$, the second 32 form a secret prefix $\mathrm{pre}$ used for nonce derivation. The public key is $A = aB$.

\textbf{Sign} message $M$:
\begin{enumerate}
  \item $r = H(\mathrm{pre} \concat M) \bmod \ell$ -- the nonce, derived deterministically from a secret prefix and the message (no RNG!).
  \item $R = rB$.
  \item $k = H(R \concat A \concat M) \bmod \ell$ -- the challenge, binding $R$, the public key, and the message.
  \item $S = (r + k\,a) \bmod \ell$.
  \item Signature is $(R, S)$, a fixed 64 bytes.
\end{enumerate}

\textbf{Verify} $(R, S)$ on $M$: compute $k = H(R\concat A\concat M)\bmod \ell$ and accept iff
\[
S\,B = R + k\,A.
\]

\begin{whybox}
\textbf{Why verification works.} By definition $S = r + ka \pmod \ell$, and scalars on this curve act mod $\ell$. So
\[
S B = (r + ka)B = rB + k(aB) = R + kA.
\]
The signer's secret scalar $a$ enters the right side only through the public $A = aB$, so the verifier checks the equation with public data alone (Figure~\ref{ch22:fig-ed25519-flow}). This is Schnorr's structure: a commitment $R$, a challenge $k$ that hashes $R$ together with the message, and a response $S$ that ties them to the secret. Hashing $R$ \emph{into} the challenge is what makes Schnorr/EdDSA cleanly provable and non-malleable, unlike (EC)DSA's bolted-together $s = k^{-1}(h + dr)$.

Try it on the small curve of Example~\ref{ch22:ex-ecdsa} ($G = (5,1)$, $n = 19$, secret $a = 7$, public $A = 7G = (0,6)$). Take nonce $r = 3$, so $R = 3G = (10,6)$, and suppose the challenge hash comes out $k = 4$. Then $S = (3 + 4\cdot 7) \bmod 19 = 31 \bmod 19 = 12$ --- no modular inverse anywhere. The verifier checks $12G = (0,11)$ against $R + 4A = 3G + 28G = 31G = 12G = (0,11)$. They match.

Why can't a forger do this without $a$? If she picks $R$ first, the hash fixes $k$, and finding $S$ with $SB = R + kA$ means solving a discrete log. If she works backward, choosing $S$ and $k$ and setting $R = SB - kA$, then $H(R \concat A \concat M)$ almost surely is not the $k$ she chose. Either order traps her.
\end{whybox}

\begin{figure}[H]
\centering
\begin{tikzpicture}[font=\scriptsize, >=stealth]
  % Seed expansion
  \node[draw=sourceblue!60, fill=sourceblue!12, rounded corners=2pt, align=center, inner sep=4pt, minimum width=2.4cm] (seed) at (0,2.6) {32-byte seed\\$\downarrow \text{SHA-512}$\\scalar $a$, prefix $\mathrm{pre}$};
  \node[draw=sourceblue!60, fill=sourceblue!12, rounded corners=2pt, align=center, inner sep=4pt, minimum width=2.4cm] (msg) at (0,0.6) {Message $M$};

  % Step 1: Deterministic nonce
  \node[draw=targetgreen!70, fill=targetgreen!12, rounded corners=2pt, align=center, inner sep=4pt, minimum width=3.1cm] (nonce) at (3.7,2.6) {\textbf{1. Deterministic Nonce}\\$r = H(\mathrm{pre}\concat M)\bmod \ell$\\$R = rB$};

  % Step 2: Challenge
  \node[draw=structurepurple!70, fill=structurepurple!12, rounded corners=2pt, align=center, inner sep=4pt, minimum width=3.1cm] (chal) at (3.7,0.6) {\textbf{2. Schnorr Challenge}\\$k = H(R\concat A\concat M)\bmod \ell$};

  % Step 3: Response
  \node[draw=quotientorange!80, fill=quotientorange!12, rounded corners=2pt, align=center, inner sep=4pt, minimum width=2.9cm] (resp) at (7.8,1.6) {\textbf{3. Linear Response}\\$S = (r + k\,a)\bmod \ell$\\(no $k^{-1}$ inversion)};

  % Step 4: Verify
  \node[draw=bobgreen!80, fill=bobgreen!15, rounded corners=2pt, align=center, inner sep=4pt, minimum width=2.9cm] (vrfy) at (11.4,1.6) {\textbf{Public Verification}\\Sig $\sigma = (R,S)$\\[2pt]Check $S\,B = R + k\,A$};

  \draw[->, thick, gray!70] (seed.east) -- (nonce.west);
  \draw[->, thick, gray!70] (msg.east) -- (chal.west);
  \draw[->, thick, gray!70] (msg.north east) -- (nonce.south west);
  \draw[->, thick, gray!70] (nonce.south) -- (chal.north);
  \draw[->, thick, gray!70] (nonce.east) -- (resp.north west);
  \draw[->, thick, gray!70] (chal.east) -- (resp.south west);
  \draw[->, thick, quotientorange!80!black] (resp.east) -- (vrfy.west);
\end{tikzpicture}
\caption{The Ed25519 signing and verification pipeline. Hashing the secret prefix $\mathrm{pre}$ with the message $M$ produces a deterministic nonce $r$ without an RNG; hashing the commitment $R=rB$ and public key $A=aB$ into the challenge $k$ prevents malleability; and the linear response $S = r + ka \bmod \ell$ eliminates modular inversion.}
\label{ch22:fig-ed25519-flow}
\end{figure}

\begin{focusbox}
What Ed25519 fixed, point by point:
\begin{itemize}
  \item \textbf{Deterministic nonce} ($r = H(\mathrm{pre}\concat M)$): the PS3 and Bitcoin disasters are structurally impossible -- there is no RNG to fail and no way to repeat $r$ across distinct messages.
  \item \textbf{No modular inversion} in signing: $S = r + ka$ has no $k^{-1}$, so no division-by-zero edge cases and easier constant-time code.
  \item \textbf{Complete addition formulas} on the Edwards curve: every point addition uses one formula with no exceptional cases, removing a class of implementation bugs and timing leaks.
  \item \textbf{Strong unforgeability and non-malleability}: the challenge hashes $R$, so the $(r, n-s)$ trick that plagues ECDSA does not apply.
  \item \textbf{Fast}: tens of thousands of signatures and verifications per second on a laptop core; batch verification is faster still.
\end{itemize}
\end{focusbox}

\begin{applicationbox}
Ed25519 is now the default host-key and user-key type in OpenSSH (\texttt{ssh-keygen -t ed25519}); it authenticates the identity keys in Signal's X3DH handshake (Chapter~27); it signs Git commits and tags, Debian and Arch package metadata, certificate-transparency logs, and many cryptocurrency and DNSSEC deployments; and it is a standard signature algorithm in TLS~1.3. If you are choosing a signature scheme for a new system in the absence of a hard interoperability or post-quantum constraint, Ed25519 is the conservative default.
\end{applicationbox}

\begin{implementbox}
A schematic of Ed25519 signing, to show the \emph{shape} of the algorithm. \textbf{For learning only -- do not use.} This sketch omits clamping details, point encoding/decoding, cofactor handling, and constant-time discipline; a real implementation (use \texttt{libsodium} or \texttt{cryptography}) gets all of those exactly right.
\begin{lstlisting}[language=Python]
import hashlib

# Assume curve helpers exist: B (base point), L (subgroup order),
# point_mul(scalar, P), point_add(P, Q), encode(P) -> bytes.
def H(b: bytes) -> int:
    return int.from_bytes(hashlib.sha512(b).digest(), "little")

def sign(A_enc: bytes, M: bytes, a: int, prefix: bytes):
    r = H(prefix + M) % L                 # deterministic nonce -- no RNG
    R = point_mul(r, B)
    R_enc = encode(R)
    k = H(R_enc + A_enc + M) % L           # challenge binds R, A, M
    S = (r + k * a) % L                    # response: NO inversion
    return R_enc + S.to_bytes(32, "little")

def verify(A, A_enc: bytes, M: bytes, sig: bytes) -> bool:
    R_enc, S = sig[:32], int.from_bytes(sig[32:], "little")
    R = decode(R_enc)
    k = H(R_enc + A_enc + M) % L
    # accept iff  S*B == R + k*A
    return point_mul(S, B) == point_add(R, point_mul(k, A))
\end{lstlisting}
The two lines worth memorizing: the nonce $r$ comes from hashing secret prefix plus message (never an RNG), and verification checks $S B = R + kA$ with no inversion in sight.
\end{implementbox}

\section{What Signatures Guarantee -- and What They Do Not}
\label{ch22:guarantees}

It is as important to know the limits as the powers. A valid signature is a precise but \emph{narrow} statement, and most real-world signature failures come from expecting it to say more than it does.

\textbf{A valid signature $\sigma$ on $m$ under $pk$ guarantees:}
\begin{itemize}
  \item \textbf{Authenticity:} someone holding the private key matching $pk$ produced $\sigma$ for this exact $m$.
  \item \textbf{Integrity:} $m$ has not changed since it was signed; one flipped bit makes verification fail.
  \item \textbf{Non-repudiation:} the signer cannot later credibly deny it (modulo key compromise).
\end{itemize}

\textbf{It does \emph{not} guarantee:}
\begin{itemize}
  \item \textbf{That $pk$ belongs to who you think.} A signature ties a message to a \emph{key}, not to a \emph{person}. Binding a key to an identity is the job of \emph{certificates} and PKI (Chapter~24). Verify a perfect signature under the wrong public key and you have authenticated an impostor.
  \item \textbf{Confidentiality.} Signatures protect integrity and origin, not secrecy. The message is typically sent in the clear alongside $\sigma$. If you need both, sign \emph{and} encrypt -- and be careful how you combine them (Chapter~16).
  \item \textbf{Freshness.} A signature on ``Pay Bob \$100'' is equally valid today and replayed next year. Signatures say nothing about \emph{when} or \emph{how many times}; defeating replay needs a nonce, timestamp, or counter \emph{inside} the signed message (the same lesson as MACs, Chapter~15).
  \item \textbf{That the signer read or meant the message.} The key endorsed these bytes; it did not promise the human understood them. ``Blind signing'' of opaque hashes has cost people real money in hardware-wallet phishing.
  \item \textbf{Anything after key compromise.} If $sk$ leaks (the PS3, the Android wallets), every past and future signature is suspect. This motivates short-lived keys, revocation (Chapter~24), and key-rotation hygiene.
\end{itemize}

\begin{trapbox}
\textbf{Sign-then-encrypt vs.\ encrypt-then-sign, and identity confusion.} Naively signing a plaintext and then encrypting it ($\Enc(\Sign(m))$) lets a recipient \emph{strip} your signature and re-encrypt the still-signed message to a third party, who is fooled into thinking you sent it to \emph{them}. Naively encrypting then signing the ciphertext lets an attacker strip your signature and re-sign the ciphertext as their own. The robust fix is to bind identities \emph{inside} the signed/encrypted content -- e.g.\ sign ``\texttt{from Alice to Bob: } $m$'' -- so the message is meaningless out of context. This ``surreptitious forwarding'' problem famously bit early PGP/email designs.
\end{trapbox}

\begin{intuitionbox}
A \emph{certificate} is a signature whose \emph{message} is ``public key $pk$ belongs to \texttt{example.com}'', signed by a certificate authority. Signatures are the building block; PKI (Chapter~24) uses them to connect a key to a name you can trust. Everything in TLS, code-signing, and the web's trust model is built by stacking signatures: a chain of ``this key vouches for that key'' all the way down to a root you decided to trust in advance.
\end{intuitionbox}

\section{Putting the Family Together}
\label{ch22:summary}

The three live signature families trade off along the same axes: underlying hard problem, key/signature size, speed, and how much rope they give a careless implementer.

\begin{center}
\renewcommand{\arraystretch}{1.3}
\begin{tabular}{@{}lllll@{}}
\toprule
\textbf{Scheme} & \textbf{Hard problem} & \textbf{Sig size (128-bit)} & \textbf{Nonce?} & \textbf{Main risk} \\
\midrule
RSA-PSS    & Factoring        & 256--384 B & none        & big keys, padding bugs \\
DSA        & DLP in $\Zns{p}$ & 64 B       & critical    & nonce reuse/bias \\
ECDSA      & EC-DLP           & 64 B       & critical    & nonce reuse, malleable \\
Ed25519    & EC-DLP (Schnorr) & 64 B       & built-in    & few; misuse-resistant \\
\bottomrule
\end{tabular}
\end{center}

The message flow is the same for all of them, and it is worth fixing in a single picture: the signer applies a private operation to a hash of the message; the verifier applies the matching public operation and checks consistency.

\begin{center}
\begin{tikzpicture}[
  font=\small,
  box/.style={draw, rounded corners, minimum height=8mm, minimum width=20mm, align=center},
  every node/.style={inner sep=3pt},
  >={Stealth[round]}
]
  % Signer side
  \node[box, fill=aliceblue!12] (m1) at (0,0) {message $m$};
  \node[box, fill=aliceblue!12] (h1) at (0,-1.6) {$\Hash(m)$};
  \node[box, fill=aliceblue!25] (sig) at (0,-3.2) {$\Sign_{sk}$\\ $\to \sigma$};
  \draw[->] (m1) -- (h1);
  \draw[->] (h1) -- (sig);
  \node[above=1mm of m1, aliceblue] {\textbf{Alice (private key $sk$)}};

  % Wire
  \node[box, fill=gray!8] (wire) at (4.2,-1.6) {send\\ $(m, \sigma)$};
  \draw[->] (m1) .. controls (2.1,0.2) and (3,-1.0) .. (wire);
  \draw[->] (sig) .. controls (2.1,-3.4) and (3,-2.2) .. (wire);
  \node[evered, below=0.5mm of wire] {Eve may read/alter};

  % Verifier side
  \node[box, fill=bobgreen!12] (h2) at (8.4,-0.6) {$\Hash(m)$};
  \node[box, fill=bobgreen!25] (vfy) at (8.4,-2.6) {$\Vrfy_{pk}(m,\sigma)$\\ check match};
  \draw[->] (wire) .. controls (6,-0.6) and (7,-0.6) .. (h2);
  \draw[->] (wire) .. controls (6,-2.6) and (7,-2.6) .. (vfy);
  \draw[->] (h2) -- (vfy);
  \node[above=1mm of h2, bobgreen] {\textbf{Bob (public key $pk$)}};
  \node[box, fill=bobgreen!10, right=6mm of vfy] (out) {accept\\ / reject};
  \draw[->] (vfy) -- (out);
\end{tikzpicture}
\end{center}

Hash to fix the message size and kill algebraic structure; apply the private operation; ship message and signature together; verify with the public key. The schemes differ only in what the ``private operation'' is -- a $d$-th power mod $N$, a discrete-log signature equation, or a Schnorr response on a curve -- and in how much care the nonce demands.

With signatures in hand, we can authenticate who is speaking across an untrusted network.
In Chapter~23, we combine ephemeral Diffie--Hellman with digital signatures to build
authenticated key exchanges with forward secrecy, so that past sessions stay secret even
if a long-term signing key is later compromised.

\section{Exercises}
\label{ch22:exercises}

\begin{exercise}[MAC vs.\ signature, in court]
Alice sends Bob a money order authenticated with a shared-key MAC. Bob takes the tag to a judge as proof Alice authorized the payment. Explain precisely why the MAC tag is \emph{not} convincing evidence, and why a digital signature would be. Which exact property is doing the work?
\end{exercise}

\begin{exercise}[Textbook RSA forgery]
With $N = 3233$, $e = 17$, $d = 2753$ (Example~\ref{ch22:ex-textbook}), an attacker has obtained Alice's textbook-RSA signatures on $m_1 = 3$ and $m_2 = 5$. Produce, with no knowledge of $d$, a valid signature on $m^* = 15 \bmod 3233$. Show the arithmetic and verify your forgery passes $\sigma^e \equiv m^* \pmod N$.
\end{exercise}

\begin{exercise}[Why hashing stops the forgery]
Repeat the previous exercise but for \emph{hash-then-sign} RSA. Explain why the multiplicative trick now fails, and state precisely what computational problem the attacker would have to solve instead (name it, and say on which object -- message or digest).
\end{exercise}

\begin{exercise}[DSA verification by hand]
Using the DSA parameters of Example~\ref{ch22:ex-dsa} ($p=23$, $q=11$, $g=4$, $y=18$), verify the signature $(r,s) = (3,10)$ on a message with digest $h = 9$ by computing $w, u_1, u_2$, and $v$ independently, and confirming $v = r$.
\end{exercise}

\begin{exercise}[Sign with ECDSA]
On the curve of Example~\ref{ch22:ex-ecdsa} ($\F_{17}$, $G=(5,1)$, $n=19$, $d=7$, $Q=(0,6)$), sign a message with digest $h = 4$ using nonce $k = 5$. You are given $5G = (9, 16)$. Compute $r$, then $s$, and write the signature. (You will need $5^{-1}\bmod 19$.)
\end{exercise}

\begin{exercise}[Recover the key from reuse]
Continuing the previous setup, suppose Alice also signed a message with digest $h' = 13$ using the \emph{same} nonce $k = 5$, producing $s'$. Without using the known value $d=7$, set up the two-signature equations, solve for $k$, then for $d$, and confirm you recover $d = 7$. (Compute $s'$ first from the signing equation, then pretend you only know $r, s, s', h, h'$.)
\end{exercise}

\begin{exercise}[Spot the bad nonce]
You are auditing a batch of ECDSA signatures from one signer. Describe the single cheap check that detects same-nonce reuse across a corpus of signatures, why it works, and what an attacker who runs the same check can then do. Relate this to how the 2013 Android-wallet thefts were carried out on the public blockchain.
\end{exercise}

\begin{exercise}[ECDSA malleability]
Show that if $(r, s)$ verifies under public key $Q$ for message digest $h$, then $(r,\, n - s)$ also verifies. (Hint: trace what happens to $w = s^{-1}$, then to $u_1, u_2$, and finally to the recovered point's $x$-coordinate.) State the ``low-$s$'' canonicalization rule that removes the ambiguity.
\end{exercise}

\begin{exercise}[Why no inversion in EdDSA]
EdDSA's response is $S = r + ka \bmod \ell$, with no $k^{-1}$, whereas ECDSA's is $s = k^{-1}(h + dr)\bmod n$. Give two concrete engineering advantages of avoiding the modular inverse (think edge cases and side channels), and explain why hashing $R$ into the challenge $k$ makes EdDSA strongly unforgeable where ECDSA is only malleably unforgeable.
\end{exercise}

\begin{exercise}[The narrow promise]
For each scenario, say whether a correctly verifying signature defends against it, and if not, name the mechanism that does: (a) an attacker flips a byte of the message in transit; (b) an attacker replays last week's signed ``unlock door'' command; (c) an attacker presents a valid signature under a key that is not actually the bank's; (d) an eavesdropper reads the signed message's contents.
\end{exercise}

\begin{exercise}[Deterministic but not random]
RFC~6979 derives $k$ deterministically from the private key and the message. Argue that this satisfies the real security requirement on the nonce even though $k$ is no longer ``random,'' and identify exactly which property of an HMAC-based derivation guarantees that two distinct messages cannot collide on the same $k$ in practice.
\end{exercise}

\begin{exercise}[Surreptitious forwarding]
Alice computes $\Enc_{pk_B}(\Sign_{sk_A}(m))$ and sends it to Bob, intending ``$m$, from Alice, for Bob.'' Show how Bob can make Carol believe Alice sent $m$ to \emph{Carol}. Then propose a concrete change to what gets signed that defeats the attack, and explain why it works.
\end{exercise}
